\chapter{La fermeture dodécaphasique de \texorpdfstring{\(\alpha\)}{alpha}}
\label{chap:alpha}
\index{alpha@\(\alpha\)}
\index{cycle dodécaphasique}

Cette section établit la fermeture dodécaphasique au moyen de l'arithmétique entière en
base mille. La construction explicite les douze récurrences, les quotients
euclidiens propagés et les conditions de frontière qui déterminent la sortie.

\section{Le sélecteur terminal des trois canaux}
\label{sec:selector-terminal-alpha-rev6}

La fermeture réunit trois générations déjà construites, mais ne les identifie pas par
leur nom décimal. À profondeur terminale, les catalogues de queues compatibles
possèdent des cardinaux
\[
 |\mathcal T_\pi|=196,\qquad
 |\mathcal T_e|=197,\qquad
 |\mathcal T_\varphi|=196.
\]
L'objet de sélection est donc
\(\mathcal T_\pi\times\mathcal T_e\times\mathcal T_\varphi\), avec
\[
 196\cdot197\cdot196=7\,567\,952
\]
triplets possibles.

La frontière de coordonnées fixée par le lecteur terminal est
\[
 q=(1,2,2,1,2,0)\in\mathbb F_3^6.
\]
La loi de canal \(\pi+e-\varphi\) possède le vecteur de signes
\[
 \sigma=(1,1,-1)=(1,1,2)\in\mathbb F_3^3
\]
et la droite duale non nulle
\[
 \langle2\sigma\rangle
 =\langle(2,2,1)\rangle.
\]
La première donnée agit sur les colonnes de la frontière ; la seconde, sur
ses lignes. Ce ne sont pas deux traits interchangeables.

\begin{theorem}[Sélecteur matriciel terminal de \(\alpha\)]
\label{thm:selector-terminal-alpha-rev6}
Dans le catalogue terminal précédent, les signatures locales de Gram, de norme, de poids et de
distance de Hamming réduisent les \(7\,567\,952\) triplets à \(331\). En exigeant
la marge de colonnes \(q\), il reste deux candidates :
\[
 (120,197,157),\qquad(129,197,148).
\]
La condition que la charge des lignes appartienne à la droite du canal
\(\langle\sigma\rangle\) sélectionne exactement la seconde. Ses blocs sont
\[
 \boxed{
 b_\pi=021020,\qquad
 b_e=001220,\qquad
 b_\varphi=100210.}
\]
\end{theorem}

\begin{proof}
Les deux matrices terminales qui survivent à la marge de colonnes sont
\[
B_0=
\begin{pmatrix}
0&2&0&2&2&0\\
0&0&1&2&2&0\\
1&0&1&0&1&0
\end{pmatrix},
\qquad
B_1=
\begin{pmatrix}
0&2&1&0&2&0\\
0&0&1&2&2&0\\
1&0&0&2&1&0
\end{pmatrix}.
\]
Leurs sommes de lignes dans \(\mathbb F_3\) sont, respectivement,
\[
 (0,2,0),\qquad(2,2,1).
\]
La première n'appartient pas à
\(\langle(1,1,2)\rangle=\{0,(1,1,2),(2,2,1)\}\) ; la seconde, si, car
\((2,2,1)=2\sigma\). Ainsi, seule \(B_1\) satisfait simultanément la
frontière de coordonnées et la charge du canal. Ses trois lignes sont les blocs
présentés et correspondent au triplet \((129,197,148)\).
\end{proof}

Le théorème n'obtient pas \(q\) à partir de la charge des lignes : \(q\) est la donnée
de coordonnées du lecteur terminal et \(\sigma\) est l'orientation du canal.
Ce qui est démontré, c'est qu'une fois déclarées les deux frontières du même objet,
la branche réelle est unique. Cette séparation empêche de transformer une signature trouvée
a posteriori en une loi prétendument dérivée.

\begin{definition}[Deux lectures de l'état terminal enrichi]
Soit
\[
 x_{\rm term}\in\mathcal X_{\rm TPK}^{\rm term,enr}
 \subseteq\mathcal X_{\rm TPK}^{[108]}.
\]
Cet état est la sortie
\[
 x_{\rm term}
 =(\mathcal U_{108}\circ\cdots\circ\mathcal U_1)(x_{\rm seed})
\]
de la chaîne complète de sélection, de transport et de mise à jour du TPK.
Le sélecteur précédent est la lecture
\[
 \pi_{\rm term}(x_{\rm term})
 =B_{\rm term}
 :=\begin{pmatrix}021020\\001220\\100210\end{pmatrix}.
\]
Le registre orienté et le lecteur des canaux de \(90^\circ\) et
\(120^\circ\) produisent, à partir du livre du même état,
\begin{equation}\label{eq:alpha-subespacio-firmado}
 \mathcal E_{108}^{90,120}\!
 \left(\mathcal R_{12}(x_{\rm term})\right)
 =(b^{90},b^{120},Q_{\rm TPK}),
\end{equation}
avec
\[
\begin{aligned}
 b^{90}={}&(-495,-116,-684,108,601,-90,
             -173,-88,313,560,-397,461),\\
 b^{120}={}&(-425,-281,-481,671,-255,30,
              475,-543,680,251,6,-128),\\
 Q_{\rm TPK}={}&6263.
\end{aligned}
\]
Pour \(r=1,2,3\), soient
\[
 B_r=\sum_{j=0}^{3}b^{120}_{r+3j},\qquad
 A_1=\frac{Q_{\rm TPK}+2B_1+B_2}{3},\qquad
 A_2=A_1-B_1,\qquad A_3=A_2-B_2,
\]
et \(d_{r,j}=b^{90}_{r+3j}\). Le lecteur harmonique \(\Pi_H\) forme les
blocs
\begin{equation}\label{eq:alpha-modulo-imagen-H12}
 \Pi_H^{(r)}
 =\bigl(A_r,d_{r,1}-d_{r,3},
              d_{r,0}+d_{r,2},d_{r,0}-d_{r,2}\bigr)
\end{equation}
et les entrelace par colonnes. La lecture signée est donc la
composition produite
\begin{equation}\label{eq:alpha-Rsgn-definido}
 \boxed{
 R_{\rm sgn}:=
 \Pi_H\circ\mathcal E_{108}^{90,120}\circ\mathcal R_{12}:
 \mathcal X_{\rm TPK}^{\rm term,enr}
 \longrightarrow
 \mathcal U_{12}^{\rm sgn}:=\operatorname{im}R_{\rm sgn}.}
\end{equation}
Ni \(U\) ni \(K\) ne sont des coordonnées ajoutées au domaine : \(U\) est la sortie
de \eqref{eq:alpha-Rsgn-definido}, et \(K\) sera son unique antécédent entier
par Hadamard.
On n'affirme pas une factorisation
\(B_{\rm term}\to U_{12}^{\rm sgn}\) : les deux sorties conservent des composantes
distinctes du même état terminal.

Pour \(N'\geq N\geq108\), les troncatures agissent uniquement sur les histoires
enrichies :
\begin{equation}\label{eq:alpha-truncamiento-firmado-preserva-KU}
 \tau_{N',N}:
 \mathcal X_{\rm TPK}^{[N'],\rm enr}
 \longrightarrow\mathcal X_{\rm TPK}^{[N],\rm enr}.
\end{equation}
Si \(R_{\rm sgn,N}\) désigne la même composition à la profondeur \(N\), la
naturalité est le carré concret
\begin{equation}\label{eq:alpha-cuadrado-naturalidad-Rsgn}
 \boxed{
 R_{\rm sgn,N}\circ\tau_{N',N}
 =R_{\rm sgn,N'}.}
\end{equation}
\end{definition}

\begin{theorem}[Construction réversible de la coordonnée signée]
\label{thm:alpha-Rsgn-construido}
La composée \eqref{eq:alpha-Rsgn-definido}, appliquée à l'état terminal
produit, donne exactement
\begin{equation}\label{eq:alpha-Rsgn-evaluacion}
 U=R_{\rm sgn}(x_{\rm term})
 =(2378,1406,2479,-452,998,-551,-668,-204,-371,-322,-28,-997).
\end{equation}
Son inversion par trois blocs de Hadamard d'ordre quatre produit ensuite
\[
 \boxed{
 K=(234,543,140,729,659,824,621,058,914,794,146,601).}
\]
La transformation est inversible ; elle satisfait le carré
\eqref{eq:alpha-cuadrado-naturalidad-Rsgn} sous les troncatures
\eqref{eq:alpha-truncamiento-firmado-preserva-KU} ; et elle est indépendante de
\(\alpha\), de CODATA et de toute valeur métrologique.
\end{theorem}

\begin{proof}
Les formules \eqref{eq:alpha-subespacio-firmado} et
\eqref{eq:alpha-modulo-imagen-H12} donnent
\[
 (B_1,B_2,B_3)=(972,-1073,101),\qquad
 (A_1,A_2,A_3)=(2378,1406,2479),
\]
et produisent les trois blocs signés
\[
\begin{aligned}
 U^{(1)}&=(2378,-452,-668,-322),\\
 U^{(2)}&=(1406,998,-204,-28),\\
 U^{(3)}&=(2479,-551,-371,-997).
\end{aligned}
\]
Leur entrelacement est \eqref{eq:alpha-Rsgn-evaluacion}. Comme
\(H_4^2=4I_4\) et
\(\det(H_4\oplus H_4\oplus H_4)=(-16)^3=-4096\), on a
\begin{equation}\label{eq:alpha-Rsgn-inversa}
 K=\mathcal H_{12}^{-1}U
 =\frac14\mathcal H_{12}U
\end{equation}
après avoir défait la permutation des orbites. En effet,
\[
 H_4U^{(1)}=(936,2916,2484,3176),\quad
 H_4U^{(2)}=(2172,2636,232,584),
\]
et
\[
 H_4U^{(3)}=(560,3296,3656,2404).
\]
La division forcée par quatre donne les orbites
\((234,729,621,794)\), \((543,659,058,146)\) et
\((140,824,914,601)\), dont l'entrelacement est le sceau déclaré.

À titre de contrôle indépendant, les opérateurs cycliques
\((D_3K)_m=K_m-K_{m+3}\) et
\((D_4K)_m=K_m-K_{m+4}\), avec \(Q(K)=\sum_mK_m\), satisfont
\begin{equation}\label{eq:alpha-extractor-90-120}
 \operatorname{rank}(D_3,D_4)=11,
 \qquad
 \operatorname{rank}(D_3,D_4,Q)=12.
\end{equation}
Ainsi \((D_3K,D_4K,Q(K))\) reconstruit un sceau unique ; de plus, le
calcul retrouve exactement
\[
 D_3K=b^{90},\qquad D_4K=b^{120},\qquad Q(K)=Q_{\rm TPK},
\]
ce qui certifie le lecteur préalable du registre, sans transformer ses sorties
en entrées. La troncature
\eqref{eq:alpha-truncamiento-firmado-preserva-KU} ne réécrit pas les cent huit premières
fenêtres du livre ; les deux membres de
\eqref{eq:alpha-cuadrado-naturalidad-Rsgn} lisent donc le même registre et
produisent \(U\). Aucune des matrices utilisées ne contient \(\alpha\) ni
une donnée expérimentale.
\end{proof}

Le calendrier observable de 108 pas est une projection strictement plus
pauvre : réinitialiser ses curseurs tous les 9 pas efface précisément les
coordonnées bidirectionnelles que le registre enrichi conserve. C'est
pourquoi le registre brut ne remplace pas l'état enrichi et n'est pas utilisé pour
reconstruire \(R_{\rm sgn}\). Le théorème de sélection précédent construit
\(B_{\rm term}\) ; le théorème \ref{thm:alpha-Rsgn-construido} construit, de
manière séparée, la seconde publication du même état.

\section{État dodécaphasique et applications réversibles}

Le théorème de cette section reçoit du
théorème~\ref{thm:alpha-Rsgn-construido} le vecteur signé
\[
\Usgn=(2378,1406,2479,-452,998,-551,-668,-204,-371,-322,-28,-997).
\]
La voie typée démontrée à partir de cette entrée est
\[
\Usgn\longleftrightarrow K,
\]
\[
K\longleftrightarrow(D_3K,D_4K,Q),
\]
\[
(P,E,\Phi,K,c_{13}=0)
\longleftrightarrow A=\alpha_{12}.
\]
L'évaluation directe ne reçoit ni \(\alpha\) ni une valeur métrologique en
entrée. Elle reçoit \(\Usgn\), les trois représentations archimédiennes et la
condition aux limites. Le
\cref{thm:selector-terminal-alpha-rev6} construit auparavant la sélection des
trois blocs terminaux ; la présente couche commence par le vecteur signé qui
les transporte à la fermeture dodécaphasique. La trace réduite d'une seule
coordonnée positionnelle n'intervient pas dans cette détermination.

\begin{remark}[Deux représentations dodécaphasiques de types distincts]
\[
U_{12}^{\rm cat}=U_6\Vert U_6
\]
est un mot périodique du catalogue, tandis que
\[
\Usgn=(2378,1406,2479,-452,998,-551,-668,-204,-371,-322,-28,-997)
\]
est un vecteur entier signé de l'espace complet des états de TPK\@. Le premier appartient à une
représentation par blocs et a une période de six. Le second conserve l'orientation, n'est pas
un mot en base mille et n'a pas cette période. Ils ne sont pas interchangeables.
\end{remark}

\section{Le vecteur canonique et les trois orbites}

Le vecteur dodécaphasique est
\[
\boxed{
K=(234,543,140,729,659,824,621,058,914,794,146,601).}
\]
Sa somme est
\[
Q(K)=6263.
\]
Le pas trois divise \(\Cdoce\) en trois orbites de quatre secteurs :
\[
\begin{aligned}
K^{(1)}&=(K_1,K_4,K_7,K_{10})=(234,729,621,794),\\
K^{(2)}&=(K_2,K_5,K_8,K_{11})=(543,659,058,146),\\
K^{(3)}&=(K_3,K_6,K_9,K_{12})=(140,824,914,601).
\end{aligned}
\]
Chaque orbite parcourt quatre secteurs séparés par \(90^\circ\). Les étiquettes
cardinales nord--est--sud--ouest apparaissent après fixation de l'origine et de l'orientation ; l'ordre
primaire est \((m,m+3,m+6,m+9)\).

\section{Première représentation : transformée de Hadamard}

Soit
\[
H_4=
\begin{pmatrix}
1&1&1&1\\
1&1&-1&-1\\
1&-1&1&-1\\
1&-1&-1&1
\end{pmatrix}.
\]
L'action sur les trois orbites donne
\[
\begin{aligned}
H_4K^{(1)}&=(2378,-452,-668,-322),\\
H_4K^{(2)}&=(1406,998,-204,-28),\\
H_4K^{(3)}&=(2479,-551,-371,-997).
\end{aligned}
\]
En entrelaçant par colonnes, nous retrouvons
\[
\mathcal H_{12}(K)=\Usgn.
\]
Comme
\[
H_4^2=4I_4,
\qquad
\det H_4=-16,
\]
la transformation globale est inversible sur \(\mathbb Q^{12}\) :
\[
K^{(r)}=\frac14H_4U^{(r)},
\qquad
\det\mathcal H_{12}=(-16)^3=-4096\ne0.
\]
Dans la représentation canonique, l'inverse produit un unique vecteur entier de
\([0,999]^{12}\).

\begin{theorem}[Structure entière de la transformation de Hadamard]
\label{thm:smith-hadamard}
La forme normale de Smith du bloc \(H_4\), notée
\(\operatorname{SNF}(H_4)\), est
\[
 \operatorname{SNF}(H_4)=\operatorname{diag}(1,2,2,4).
\]
Par conséquent, pour la somme directe des trois blocs,
\[
 \operatorname{coker}\mathcal H_{12}
 \cong(\mathbb Z/2\mathbb Z)^6\oplus(\mathbb Z/4\mathbb Z)^3,
\]
et l'image de \(\mathbb Z^{12}\) a pour indice \(4096\). Un vecteur
\(u\in\mathbb Z^4\) appartient à \(H_4\mathbb Z^4\) exactement lorsque
\[
 H_4u\equiv0\pmod 4.
\]
\end{theorem}

\begin{proof}
Le plus grand commun diviseur des entrées de \(H_4\) est un. Les mineurs
d'ordre deux sont pairs et l'un a pour valeur absolue deux ; les mineurs d'ordre
trois sont des multiples de quatre et l'un a pour valeur absolue quatre ;
enfin, \(\lvert\det H_4\rvert=16\). Les diviseurs déterminantaux sont
 donc \(1,2,4,16\), d'où l'on obtient les facteurs invariants
\(1,2,2,4\). La somme directe de trois blocs répète ces facteurs et a pour
indice \(16^3=4096\). Puisque \(H_4^{-1}=H_4/4\), la dernière
congruence équivaut à l'intégralité de l'antécédent.
\end{proof}

\section{Transporteur cyclotomique relatif entre les secteurs trois et quatre}
\label{sec:transportador-ciclotomico-a34-rev8}

Soit \(\mathsf R\) la rotation régulière du cycle \(C_{12}\) sur
\(\mathbb Q^{12}\), et soient
\[
 W_d^{\rm cyc}:=\ker\Phi_d(\mathsf R)
\]
ses secteurs cyclotomiques. Les projecteurs rationnels sur les secteurs
d'ordres trois et quatre peuvent s'écrire
\[
 \Pi_3=
 \frac14(I+\mathsf R^3+\mathsf R^6+\mathsf R^9)
 -\frac1{12}\sum_{k=0}^{11}\mathsf R^k,
\]
\[
 \Pi_4=
 \frac16(I-\mathsf R^2+\mathsf R^4-\mathsf R^6
              +\mathsf R^8-\mathsf R^{10}).
\]
Pour fixer l'ordre des coordonnées, soit
\[
 \mathsf P_{\rm orb}(z_1,\ldots,z_{12})
 =\bigl((z_1,z_4,z_7,z_{10}),
        (z_2,z_5,z_8,z_{11}),
        (z_3,z_6,z_9,z_{12})\bigr).
\]
La transformée globale utilisée ci-dessus est, dans l'ordre naturel de \(C_{12}\),
\[
 \mathcal H_{12}
 :=\mathsf P_{\rm orb}^{-1}
   (H_4\oplus H_4\oplus H_4)\mathsf P_{\rm orb}.
\]
Autrement dit, \(H_4\oplus H_4\oplus H_4\) agit dans l'ordre des trois
orbites de pas trois, et non sur trois quadruplets contigus de l'ordre naturel.
Nous définissons alors le transporteur relatif
\[
 A_{34}:=\left.\Pi_4\mathcal H_{12}\Pi_3
          \right|_{W_3^{\rm cyc}}:
 W_3^{\rm cyc}\longrightarrow W_4^{\rm cyc}.
\]

\begin{theorem}[Transporteur cyclotomique dodécaphasique]
\label{thm:transportador-ciclotomico-a34-rev8}
Dans les bases
\[
 u_1=(1,-1,0)^{\times4},\qquad
 u_2=(0,1,-1)^{\times4}
\]
de \(W_3^{\rm cyc}\), et
\[
 v_1=(1,0,-1,0)^{\times3},\qquad
 v_2=(0,1,0,-1)^{\times3}
\]
de \(W_4^{\rm cyc}\), on a
\[
 A_{34}u_1=\frac23(v_1-v_2),\qquad
 A_{34}u_2=\frac23(v_1+v_2),
\]
et, par conséquent,
\[
 \boxed{\det A_{34}=\frac89}.
\]
En particulier, \(A_{34}\) est un isomorphisme rationnel entre les deux secteurs.
\end{theorem}

\begin{proof}
On substitue les quatre suites périodiques précédentes dans les expressions
de \(\Pi_3,\Pi_4\) et dans les trois blocs de Hadamard. La matrice de \(A_{34}\)
dans les bases déclarées est
\[
 \begin{pmatrix}
  2/3&2/3\\
 -2/3&2/3
 \end{pmatrix},
\]
dont le déterminant est \(8/9\).
\end{proof}

L'involution rationnelle qui échange \(W_3^{\rm cyc}\) et
\(W_4^{\rm cyc}\) au moyen de \(A_{34}\) et de \(A_{34}^{-1}\), et fixe
\(W_{12}^{\rm cyc}\), transporte la somme
\(W_4^{\rm cyc}\oplus W_{12}^{\rm cyc}\) vers
\(W_3^{\rm cyc}\oplus W_{12}^{\rm cyc}\). Dans cette dernière somme, le polynôme
cyclotomique de la rotation est
\[
 \Phi_3(x)\Phi_{12}(x),
\]
qui coïncide avec le polynôme du Coxeter de type \(E_6\). C'est une signature
spectrale exacte qui prépare le corridor géométrique ultérieur.

\begin{statusbox}[title={Transport rationnel et continuation entière}]
L'opérateur \(A_{34}\) réalise le transport cyclotomique rationnel entre les
secteurs d'ordres trois et quatre. En fixant \(W_{12}^{\rm cyc}\), il situe la
dodécaphase dans le secteur spectral exact
\(\Phi_3\Phi_{12}\) de type \(E_6\). La continuation entière est construite
au moyen de deux sources ultérieures et compatibles : le
\cref{thm:entrelazador-dodecafase-witt} fournit l'isométrie
\(M_{12}\)-équivariante vers le module de Witt, et le
\cref{thm:identificacion-isometrica-tres-proyectores} enchaîne les trois
projecteurs d'incidence. Sur cette réalisation entière, le
\cref{p12-83-c20-s6:thm:grafo-27-rectas-nuevo} construit le graphe des
vingt-sept droites et le
\cref{p12-83-c20-s6:chap:holonomia-schlafli-e6-nuevo} identifie son action de
Weyl. Ainsi, \(A_{34}\) fixe l'étape rationnelle de la généalogie et les
opérateurs d'entrelacement de Witt--Schläfli réalisent son prolongement réticulaire et
incidenciel.
\end{statusbox}

La séparation des signes est
\[
U_{12}^{+}
=(2378,1406,2479,0,998,0,0,0,0,0,0,0),
\]
\[
U_{12}^{-}
=(0,0,0,452,0,551,668,204,371,322,28,997),
\]
et
\[
\Usgn=U_{12}^{+}-U_{12}^{-}.
\]

\section{Seconde représentation : différences transversales}

Avec des indices dans \(\{1,\ldots,12\}\) et la convention cyclique
\[
 K_{m+r}:=K_{1+((m+r-1)\bmod12)},
\]
nous définissons
\[
(D_3K)_m=K_m-K_{m+3},
\qquad
(D_4K)_m=K_m-K_{m+4}.
\]
L'extracteur transversal dodécaphasique est l'opérateur
\[
 \mathcal E_{\mathrm{dod}}
 :=(D_3,D_4,Q).
\]
L'indice ``dod'' déclare son domaine. Les décalages trois et quatre
correspondent à des séparations angulaires de \(90^\circ\) et \(120^\circ\) dans le
cycle orienté \(C_{12}\), mais \(\mathcal E_{\mathrm{dod}}\) n'est ni l'opérateur
d'incidence hexade--octade ni le noyau angulaire utilisé ultérieurement dans
les tours.
Les données sont
\[
\begin{aligned}
D_3K={}&(-495,-116,-684,108,601,-90,\\
       &\qquad-173,-88,313,560,-397,461),\\[1mm]
D_4K={}&(-425,-281,-481,671,-255,30,\\
       &\qquad475,-543,680,251,6,-128),\\[1mm]
Q={}&6263.
\end{aligned}
\]
Si \(D_3v=D_4v=0\), le vecteur est constant aux pas trois et quatre. Comme
\(\gcd(3,4)=1\), ces pas relient tout le cycle et
\[
\ker D_3\cap\ker D_4=\langle\mathbf1\rangle.
\]
Par conséquent
\[
\operatorname{rank}(D_3,D_4)=11.
\]
La charge totale élimine la translation constante, car
\(Q(\mathbf1)=12\ne0\). Donc
\[
\boxed{\operatorname{rank}(D_3,D_4,Q)=12,}
\]
et
\[
K\longleftrightarrow(D_3K,D_4K,Q)
\]
est un système de coordonnées exact.

\begin{theorem}[Forme normale entière du système transversal]
\label{thm:smith-dodecafase}
Soit
\[
 \mathcal D=(D_3,D_4,Q):
 \mathbb Z^{12}\longrightarrow\mathbb Z^{25}.
\]
Ses facteurs invariants non nuls sont
\[
 \operatorname{SNF}(\mathcal D)
 =\operatorname{diag}(\underbrace{1,\ldots,1}_{11},12),
\]
et, par conséquent,
\[
 \operatorname{coker}\mathcal D
 \cong\mathbb Z^{13}\oplus\mathbb Z/12\mathbb Z.
\]
\end{theorem}

\begin{proof}
Les lignes de \((D_3,D_4)\) sont des incidences orientées du graphe
\[
 \operatorname{Cay}
 \bigl(\mathbb Z/12\mathbb Z,\{\pm3,\pm4\}\bigr).
\]
Comme les pas trois et quatre engendrent \(\mathbb Z/12\mathbb Z\), le graphe
est connexe.
Les incidences d'un arbre couvrant fournissent onze mineurs réduits
unimodulaires ; les onze premiers facteurs invariants sont donc égaux à
un.

Tout mineur d'ordre douze contient la ligne \(Q\), puisque le bloc
d'incidences est de rang onze. En ajoutant les onze premières colonnes à la
dernière, opération unimodulaire, la dernière colonne s'annule dans les lignes
d'incidence et vaut douze dans la ligne \(Q\). Tous les mineurs maximaux sont
divisibles par douze. Le choix des onze incidences d'un arbre donne un
mineur restant de valeur absolue un et, par conséquent, un déterminant maximal
de valeur absolue douze. Le dernier facteur invariant est exactement douze.
\end{proof}

La composante \(\mathbb Z/12\mathbb Z\) enregistre la condition de
compatibilité entière de la charge totale avec le potentiel reconstruit.
En particulier, l'unicité rationnelle prouvée par le rang est affinée ici en
un énoncé sur le réseau entier et son unique obstruction torsionnelle.

L'observable signé peut aussi être reconstruit à partir de ces coordonnées.
Pour \(r=1,2,3\), soit
\[
B_r=\sum_{j=0}^{3}(D_4K)_{r+3j}.
\]
On obtient
\[
(B_1,B_2,B_3)=(972,-1073,101).
\]
Les sommes orbitales sont
\[
A_1=\frac{Q+2B_1+B_2}{3}=2378,
\quad
A_2=A_1-B_1=1406,
\quad
A_3=A_2-B_2=2479.
\]
Si
\(d_{r,j}=(D_3K)_{r+3j}\) pour \(j=0,1,2,3\), avec la même convention
cyclique, le bloc signé de l'orbite \(r\) est
\[
(A_r,d_{r,1}-d_{r,3},d_{r,0}+d_{r,2},d_{r,0}-d_{r,2}),
\]
exactement le bloc de Hadamard correspondant.

\begin{remark}[Codomaines de \(\mathcal H_{12}\) et
\(\mathcal E_{\mathrm{dod}}\)]
L'écriture correcte est
\[
\mathcal H_{12}(K)=\Usgn,
\qquad
\mathcal E_{\mathrm{dod}}(K)=(D_3K,D_4K,Q).
\]
Les deux représentations sont compatibles car elles reconstruisent le même vecteur,
mais aucun diagramme entre leurs codomaines n'a été défini ici. Leurs formules et
codomaines sont distincts. Partager les pas \(3/4\) ou les indices \(90/120\) n'autorise pas
à les identifier.
\end{remark}

\section{Les trois coordonnées archimédiennes}

Les vecteurs de douze triades sont
\[
\begin{aligned}
P={}&(141,592,653,589,793,238,462,643,383,279,502,884),\\
E={}&(718,281,828,459,045,235,360,287,471,352,662,497),\\
\Phi={}&(618,033,988,749,894,848,204,586,834,365,638,117).
\end{aligned}
\]
Sous forme concaténée :
\[
\begin{aligned}
\iota(P)&=141592653589793238462643383279502884,\\
\iota(E)&=718281828459045235360287471352662497,\\
\iota(\Phi)&=618033988749894848204586834365638117.
\end{aligned}
\]
\(P\) est la coordonnée décimale partagée par les cinq régions de
\(\pi\) ; le calcul n'affirme pas que ces régions sont identiques.

\section{Récurrence de retenue en base mille}

Nous fixons
\[
c_{13}=0
\]
et calculons de droite à gauche, \(m=12,11,\ldots,1\) :
\[
s_m=P_m+E_m-\Phi_m-K_m+c_{m+1},
\]
\[
A_m=s_m\bmod1000,
\qquad
c_m=\frac{s_m-A_m}{1000},
\qquad0\le A_m<1000.
\]
Le tableau est imprimé dans l'ordre croissant des indices, mais chaque retenue provient de la ligne
inférieure.

\begingroup
\scriptsize
\setlength{\tabcolsep}{3.3pt}
\begin{longtable}{@{}r|rrrr|rr|rr@{}}
\caption{Fermeture dodécaphasique complète en base mille.}\label{tab:acarreo-alpha}\\
\toprule
\(m\)&\(P_m\)&\(E_m\)&\(\Phi_m\)&\(K_m\)&\(c_{m+1}\)&\(s_m\)&\(A_m\)&\(c_m\)\\
\midrule
\endfirsthead
\toprule
\(m\)&\(P_m\)&\(E_m\)&\(\Phi_m\)&\(K_m\)&\(c_{m+1}\)&\(s_m\)&\(A_m\)&\(c_m\)\\
\midrule
\endhead
1&141&718&618&234&0&7&007&0\\
2&592&281&033&543&0&297&297&0\\
3&653&828&988&140&\(-1\)&352&352&0\\
4&589&459&749&729&\(-1\)&\(-431\)&569&\(-1\)\\
5&793&045&894&659&\(-2\)&\(-717\)&283&\(-1\)\\
6&238&235&848&824&\(-1\)&\(-1200\)&800&\(-2\)\\
7&462&360&204&621&0&\(-3\)&997&\(-1\)\\
8&643&287&586&058&\(-1\)&285&285&0\\
9&383&471&834&914&\(-1\)&\(-895\)&105&\(-1\)\\
10&279&352&365&794&0&\(-528\)&472&\(-1\)\\
11&502&662&638&146&0&380&380&0\\
12&884&497&117&601&0&663&663&0\\
\bottomrule
\end{longtable}
\endgroup

Les retenues sont
\[
(c_1,\ldots,c_{13})
=(0,0,0,-1,-1,-2,-1,0,-1,-1,0,0,0).
\]
Les valeurs négatives de \(s_m\) sont des emprunts ordinaires en base mille. La
division euclidienne fixe de manière unique le reste et le quotient.

La sortie est
\[
\boxed{
A=(007,297,352,569,283,800,997,285,105,472,380,663),}
\]
et donc
\[
\boxed{
\begin{aligned}
\alpha_{12}
&:=\frac{7297352569283800997285105472380663}{10^{36}},\\
&=0.007297352569283800997285105472380663.
\end{aligned}}
\]
Ce rationnel est la fenêtre dodécaphasique initiale de la section compatible
engendrée par le transport TPK. Si
\(\rho_{N',N}:C_{N'}^\alpha\to C_N^\alpha\) sont ses troncatures, la
sortie complète est
\[
 \boxed{
 \alpha_{\mathrm{HMT}}
 :=\varprojlim_N\alpha_N,\qquad
 \rho_{N',N}(\alpha_{N'})=\alpha_N,\qquad
 \alpha_{36}=\alpha_{12}.}
\]
Le retour de phase prolonge la mémoire et resserre les cylindres
survivants à une profondeur arbitraire ; il ne réinitialise pas l'état et n'achève pas la
construction à trente-six chiffres.

\section{L'identité affine complète}

Coordonnée par coordonnée,
\[
\boxed{
P_m+E_m-\Phi_m-K_m+c_{m+1}
=A_m+1000c_m.}
\]
Si
\[
C=(c_1,\ldots,c_{12}),
\qquad
C^+=(c_2,\ldots,c_{13}),
\]
la forme vectorielle correcte est
\[
\boxed{P+E-\Phi-K+C^+-1000C=A.}
\]
Pour éviter d'identifier les queues finies aux constantes complètes,
nous définissons
\[
 \tau_{36}(x)=10^{-36}\left\lfloor10^{36}
 \bigl(x-\lfloor x\rfloor\bigr)\right\rfloor,
 \qquad
 \kappa(K)=10^{-36}\iota(K).
\]
La notation décimale exacte de l'identité concaténée est alors
\[
 \tau_{36}(\pi)+\tau_{36}(e)-\tau_{36}(\varphi)-\kappa(K)
 =\alpha_{12}.
\]
Sans la troncature explicite, le plongement \(\kappa\) et le cocycle
\(C^+-1000C\), l'égalité en coordonnées serait fausse.

On définit
\[
\iota(v_1,\ldots,v_{12})
=\sum_{m=1}^{12}v_m1000^{12-m}.
\]
En multipliant chaque récurrence par \(1000^{12-m}\) et en sommant, les retenues
se télescopent :
\[
\iota(P)+\iota(E)-\iota(\Phi)-\iota(K)-\iota(A)
=1000^{12}c_1-c_{13}.
\]
Comme \(c_1=c_{13}=0\) :
\[
\boxed{
\iota(P)+\iota(E)-\iota(\Phi)-\iota(K)=\iota(A).}
\]
Les deux entiers restants sont
\[
\iota(K)=234543140729659824621058914794146601,
\]
\[
\iota(A)=007297352569283800997285105472380663.
\]

\section{Inversion et équivalence des trois représentations}

La récurrence s'inverse par
\[
K_m=P_m+E_m-\Phi_m+c_{m+1}-A_m-1000c_m,
\]
et la représentation entière concaténée par
\[
\iota(K)=\iota(P)+\iota(E)-\iota(\Phi)-\iota(A).
\]
Le diagramme complet est
\[
\boxed{
K\longleftrightarrow\Usgn,
\qquad
K\longleftrightarrow(D_3K,D_4K,Q),
\qquad
K\longleftrightarrow A\mid(P,E,\Phi,c_{13}=0).}
\]
Les deux premières sont des représentations linéaires ; la troisième, une application
affine réversible dans la fibre
conditionnée par \(P,E,\Phi\) et la frontière.

\begin{remark}[Commutativité de la fermeture]
L'application directe obtient \(K\) à partir de la représentation signée et produit
\(A\). La formule inverse démontre que la sortie et la frontière retrouvent le
même vecteur sans liberté résiduelle. L'inversibilité ne constitue pas une
preuve statistique indépendante ; la commutativité des trois
représentations est une propriété algébrique de la fermeture.
\end{remark}

L'unicité s'exprime à trois niveaux de la même construction :
\begin{enumerate}
  \item une fois fixés le vecteur signé, \(P,E,\Phi\) et la frontière, le vecteur et la
  sortie sont uniques ;
  \item les représentations sont réversibles et commutatives ;
  \item la généalogie complète
  \(\mathrm{APP}\to\mathrm{TRIT}\to\mathrm{TPK}\) produit ces primitives
  avant la fermeture locale, et ses troncatures compatibles prolongent
  l'identité à la section coinductive \(\alpha_{\mathrm{HMT}}\).
\end{enumerate}

\section{Condition aux limites dodécaphasique}

La première fenêtre dodécaphasique utilise \(c_{13}=0\) et se termine en
\[
\ldots|380|663.
\]
Cette frontière fait partie du domaine de l'application affine réversible. Le
prolongement n'introduit pas d'autre système : il itère le même transport gradué
du TPK sur
\[
 w_6\to w_{12}\to w_{18}\to w_{24}\to w_{30}\to R_{36}\to G_9
\]
et construit les fenêtres successives avec retour de phase, avancée de la mémoire et
retenue terminale compatible. Le
\cref{sec:l9s102:alpha-dos-vias} démontre que les carrés de troncature
de cette voie et du lecteur analytique commutent et que leurs cylindres
survivants possèdent une intersection réduite à un singleton.

\subsection{Classe entière de la retenue}

La condition terminale peut être comparée au problème périodique. Soit
\(S\) le décalage cyclique sur \(\mathbb Z^{12}\) et, pour une base
entière \(b\geq2\), définissons
\[
 \partial_b=S-bI.
\]
L'identité affine périodique prend la forme
\[
 A=P+E-\Phi-K+\partial_bC.
\]
Pour tout \(h\in\mathbb Z^{12}\), la transformation
\[
 K\longmapsto K+\partial_bh,
 \qquad
 C\longmapsto C+h
\]
conserve \(A\). Ainsi, le choix des retenues est un choix de
représentant au sein d'une classe entière.

\begin{proposition}[Quotient périodique de la retenue]
\label{prop:cociente-periodico-acarreo}
Sur un cycle de longueur douze,
\[
 \operatorname{coker}(S-bI)\cong
 \mathbb Z/(b^{12}-1)\mathbb Z.
\]
La concaténation en base \(b\), prise modulo \(b^{12}-1\), détermine la
classe périodique.
\end{proposition}

\begin{proof}
Le module cyclique s'identifie à
\(\mathbb Z[t]/(t^{12}-1)\). Imposer la relation \(t=b\) produit
\(\mathbb Z/(b^{12}-1)\mathbb Z\). L'évaluation du polynôme des
coefficients est la concaténation en base \(b\), avec l'orientation fixée
par l'ordre des douze secteurs.
\end{proof}

La fermeture de \(\alpha\) n'utilise pas le quotient périodique : elle utilise une chaîne
orientée avec une extrémité droite et exige \(c_{13}=0\). La division euclidienne
sélectionne alors un représentant unique. Cette différence explique
pourquoi la frontière ouverte fait partie de la définition de l'application et ne
peut pas être remplacée par une identification cyclique des retenues.

\section{Comparaison avec la référence métrologique}

Après la génération, l'arithmétique entière de la première fenêtre
vérifie la reconnaissance
\[
\boxed{
\alpha_{12}
=10^{-36}
\left\lfloor\frac{10^{36}}{137.035999084}\right\rfloor.}
\]
Son inverse est
\[
\alpha_{12}^{-1}
=137.03599908400000000000000000000001066588\ldots,
\]
et reproduit exactement la chaîne centrale publiée dans l'ajustement de 2018 du
Committee on Data for Science and Technology (CODATA) \cite{Tiesinga2021CODATA} :
\[
\boxed{137.035999084.}
\]
Il ne s'agit pas d'une coïncidence en virgule flottante. Les trente-six chiffres de
\(\alpha_{12}\) constituent la troncature exacte de l'inverse de la
chaîne centrale imprimée. Cette vérification finie ne définit pas
\(\alpha_{\mathrm{HMT}}\) et ne limite pas sa profondeur : la section coinductive
déjà produite détermine toutes ses fenêtres compatibles, tandis que CODATA ne
reconnaît qu'ensuite le préfixe que sa résolution expérimentale permet
de confronter.

La reconnaissance métrologique ne fait pas partie de la récurrence. Le résultat
mathématique est l'identité exacte des deux voies
\[
 \alpha_A=\alpha_B=\alpha_{\mathrm{HMT}},
\]
dont les fenêtres dodécaphasiques sont déterminées par le vecteur signé, les
trois mots engendrés et les frontières compatibles. Le tableau de CODATA
ne compare qu'ensuite une sortie déjà fixée.

La comparaison avec des ajustements métrologiques ultérieurs est un contrôle externe
séparé. La sortie HMT reste figée et peut donc être confrontée sans
modifier le générateur.

\section{Indépendance des coordonnées et rigidité}

Les tests de la fermeture distinguent les modifications profondes des perturbations
locales :
\begin{center}
\small
\begin{tabularx}{\textwidth}{@{}lYY@{}}
\toprule
Variante & Résultat & Ce qu'il contrôle \\
\midrule
protocole canonique droite--gauche & igualdad exacta & référence \\
utiliser \(U_6\Vert U_6\) comme \(K\) & échoue macroscopiquement & distinction des types \\
rotar \(K\) & rompt la fermeture & origine dodécaphasique \\
reflejar \(K\) & inverse le bord & orientation \\
retenue gauche--droite & change le mot & sens du transport \\
perturber \(K_i\) d'une unité & le mot exact échoue & rigidez profunda \\
\bottomrule
\end{tabularx}
\end{center}
Une perturbation locale peut ne modifier qu'un chiffre profond ; elle reste une
réfutation de l'égalité exacte. Les variantes structurelles produisent des échecs
plus importants et montrent que la fermeture n'est pas invariante sous des ordres arbitraires.

\begin{remark}[Structure de la fermeture dodécaphasique]
La fermeture de \(\alpha\) comprend un cycle orienté, deux représentations
exactes du vecteur \(K\), douze entrées typées, une récurrence de droite à
gauche, treize retenues, une condition aux limites et une identité entière
télescopique. Ces données déterminent la sortie de manière reproductible.
\end{remark}
